\documentclass[12pt,a4paper]{article}

\usepackage[margin=2.3cm]{geometry}
\usepackage{fontspec}
\usepackage{xepersian}
\usepackage{longtable}
\usepackage{tabularx}
\usepackage{booktabs}
\usepackage{array}
\usepackage{xcolor}
\usepackage{hyperref}
\usepackage{enumitem}
\usepackage{fancyhdr}
\usepackage{titlesec}
\usepackage{fancyvrb}

\providecommand{\UseMathForPositioningText}{}

\IfFontExistsTF{Vazirmatn}{
  \settextfont{Vazirmatn}
}{
  \IfFontExistsTF{Tahoma}{
    \settextfont{Tahoma}
  }{
    \settextfont{Arial}
  }
}

\IfFontExistsTF{Consolas}{
  \setmonofont{Consolas}
}{
  \setmonofont{Courier New}
}

\IfFontExistsTF{Arial}{
  \setlatintextfont{Arial}
}{
  \setlatintextfont{Times New Roman}
}

\hypersetup{
  colorlinks=true,
  linkcolor=blue,
  urlcolor=blue,
  pdftitle={AI Chat SaaS System Architecture},
  pdfauthor={Mobin Yousefi}
}

\definecolor{HeaderBlue}{HTML}{1E3A8A}
\definecolor{SoftGray}{HTML}{F8FAFC}
\definecolor{BorderGray}{HTML}{CBD5E1}

\setlength{\parindent}{0pt}
\setlength{\parskip}{7pt}
\setlength{\headheight}{15pt}
\emergencystretch=3em
\sloppy
\renewcommand{\arraystretch}{1.45}
\setlist[itemize]{rightmargin=1.2cm,itemsep=3pt,topsep=4pt}
\setlist[enumerate]{rightmargin=1.2cm,itemsep=3pt,topsep=4pt}

\titleformat{\section}{\Large\bfseries\color{HeaderBlue}}{\thesection}{0.6em}{}
\titleformat{\subsection}{\large\bfseries}{\thesubsection}{0.6em}{}
\titleformat{\subsubsection}{\normalsize\bfseries}{\thesubsubsection}{0.6em}{}

\pagestyle{fancy}
\fancyhf{}
\rhead{مستند معماری سیستم AI Chat SaaS}
\lhead{نسخه \en{1.0}}
\cfoot{\thepage}

\newcommand{\en}[1]{\lr{#1}}
\newcommand{\code}[1]{\lr{\texttt{#1}}}

\newenvironment{diagramblock}
{%
  \VerbatimEnvironment
  \begin{latin}
  \begin{Verbatim}[
    fontsize=\small,
    frame=single,
    rulecolor=\color{BorderGray},
    baselinestretch=1,
    xleftmargin=0.2cm,
    xrightmargin=0.2cm
  ]%
}
{%
  \end{Verbatim}
  \end{latin}
}

\begin{document}

\begin{titlepage}
\centering
\vspace*{2.2cm}
{\Huge\bfseries مستند معماری سیستم \en{AI Chat SaaS}\par}
\vspace{0.7cm}
{\Large\bfseries \en{System Architecture Documentation}\par}
\vspace{1.3cm}
{\large معماری سامانه چندمستاجری چت پشتیبانی، پنل مدیریتی، ویجت سایت و دستیار هوشمند مبتنی بر دانش\par}
\vfill
\begin{tabular}{rl}
\textbf{نام پروژه:} & \en{AI Chat SaaS} \\
\textbf{نوع سند:} & مستند معماری سیستم \\
\textbf{نسخه سند:} & \en{1.0} \\
\textbf{تاریخ:} & ۱۴ جولای ۲۰۲۶ \\
\textbf{نویسنده/تهیه‌کننده:} & مبین یوسفی \\
\textbf{مسیر سند:} & \code{doc/LaTeX/04-System-Architecture/04-System-Architecture.tex} \\
\end{tabular}
\vfill
{\small این سند برای ارائه معماری فنی، اجزای سیستم، جریان ارتباطی، مدل استقرار، مدل داده و مرزهای امنیتی سامانه تهیه شده است.}
\end{titlepage}

\tableofcontents
\newpage

\section{هدف و دامنه سند}
این سند معماری سیستم \en{AI Chat SaaS} را تشریح می‌کند. تمرکز سند بر نمایش ساختار واقعی محصول، اجزای نرم‌افزاری، ارتباط بین لایه‌ها، جریان‌های اصلی، مدل داده، استقرار، شبکه، امنیت و تصمیم‌های معماری است.

مطابق راهنمای \code{Rahnama.md}، این سند شامل موارد زیر است:
\begin{itemize}
  \item \en{Architecture Diagram}
  \item \en{Context Diagram}
  \item \en{Component Diagram}
  \item \en{Deployment Diagram}
  \item \en{Sequence Diagram}
  \item \en{Class/Domain Diagram}
  \item \en{Microservice Diagram} به‌صورت نقشه ماژول‌های سرویس‌مانند در معماری فعلی
  \item \en{Network Diagram}
  \item \en{Database Diagram}
\end{itemize}

منابع فنی این سند شامل \code{ArchitectureForMobin.md}، \code{ReportForSale.md}، endpointهای \code{backend/api}، helperهای \code{backend/includes}، migration مربوط به AI و ساختار frontend و widget است.

\section{خلاصه معماری}
\en{AI Chat SaaS} یک سامانه B2B SaaS چندمستاجری است که از سه سطح اصلی تشکیل شده است:
\begin{enumerate}
  \item \textbf{سطح عمومی سایت مشتری:} ویجت جاوااسکریپت که روی سایت مشتری نصب می‌شود و با API عمومی widget ارتباط دارد.
  \item \textbf{سطح عملیات پشتیبانی:} پنل Next.js برای سوپرادمین، ادمین مشتری و agentها.
  \item \textbf{سطح بک‌اند و داده:} APIهای PHP، پایگاه داده MySQL، فضای upload و موتور دانش/AI.
\end{enumerate}

بک‌اند در نسخه فعلی به‌صورت endpointهای مستقل PHP طراحی شده است. هر endpoint مسئول یک عملیات مشخص است و قابلیت‌های مشترک مانند احراز هویت، پاسخ JSON، CORS، کنترل دسترسی سایت، rate limit، upload و موتور AI از طریق فایل‌های \code{backend/includes} استفاده می‌شوند.

\newpage
\section{نمودار معماری کلان \en{(Architecture Diagram)}}
\textbf{نمودار ۱: نمای معماری کلان سامانه}

\begin{diagramblock}
Website Visitor       Admin / Agent        Super Admin
      |                    |                    |
      v                    v                    v
+-------------+      +-------------------------------+
| Chat Widget |      | Next.js Management Panel      |
| JS + Shadow |      | Super Admin / Customer / Agent |
+-------------+      +-------------------------------+
      |                    |
      v                    v
+-------------+      +-------------------------------+
| Widget API  |      | Panel APIs                    |
| site_key    |      | auth / agent / customer / SA  |
+-------------+      +-------------------------------+
      \                    /
       \                  /
        v                v
   +---------------------------------------------+
   | PHP Backend Core                            |
   | auth, jwt, cors, site access, plan limits,  |
   | rate limit, upload, response, ai helpers    |
   +---------------------------------------------+
        |                |                 |
        v                v                 v
 +-------------+   +----------------+   +----------------+
 | MySQL DB    |   | AI Knowledge   |   | Upload Storage |
 | SaaS data   |   | crawler/search |   | attachments    |
 +-------------+   +----------------+   +----------------+
\end{diagramblock}

\subsection{توضیح معماری کلان}
در این معماری، ویجت و پنل دو client مستقل هستند. ویجت از طریق \code{site\_key} و کنترل origin به API عمومی وصل می‌شود. پنل مدیریتی با JWT و نقش کاربر به endpointهای محافظت‌شده متصل می‌شود. بک‌اند، مرز اصلی اعمال قوانین دسترسی، محدودیت پلن، امنیت فایل، ثبت مکالمه و اجرای منطق AI است.

\newpage
\section{نمودار زمینه و مرز سیستم \en{(Context Diagram)}}
\textbf{نمودار ۲: Context Diagram و مرز سامانه}

\begin{diagramblock}
                   +-------------------------+
                   | Customer Website Pages  |
                   | HTML used for crawling  |
                   +-----------+-------------+
                               |
                               v
+-------------+       +-----------------------+       +-------------+
| Visitor     | ----> | AI Chat SaaS System   | ----> | MySQL DB    |
| chat user   |       | widget + panel + API  |       | data store  |
+-------------+       +-----------------------+       +-------------+
      ^                         ^       |
      |                         |       v
+-------------+       +---------+---+   +----------------+
| Customer    | ----> | Platform   |   | File Storage   |
| Admin/Agent |       | SuperAdmin |   | uploads        |
+-------------+       +-------------+   +----------------+
\end{diagramblock}

مرز سیستم در این نمودار شامل frontend، widget، backend، دیتابیس و storage است. سایت مشتری و مرورگرهای کاربران بیرون از مرز سیستم قرار دارند، اما از طریق script نصب‌شده و APIهای کنترل‌شده با سامانه تعامل می‌کنند.

\newpage
\section{نمودار اجزا \en{(Component Diagram)}}
\textbf{نمودار ۳: Component Diagram}

\begin{diagramblock}
+----------------------+        +--------------------------+
| Frontend / Next.js   |        | Widget / JavaScript      |
| app pages, api.ts    |        | Shadow DOM, polling      |
+----------+-----------+        +------------+-------------+
           |                                 |
           v                                 v
+----------------------+        +--------------------------+
| Protected Panel APIs |        | Public Widget APIs       |
| auth, agent,         |        | config, visitor,         |
| customer, superadmin |        | conversation, ai-reply   |
+----------+-----------+        +------------+-------------+
           |                                 |
           +---------------+-----------------+
                           v
             +-----------------------------+
             | Shared Backend Includes     |
             | auth, jwt, cors, response,  |
             | site-access, plan-limits,   |
             | upload, rate-limit, ai      |
             +--------------+--------------+
                            |
        +-------------------+-------------------+
        v                   v                   v
+---------------+   +----------------+   +----------------+
| MySQL Tables  |   | AI Engine      |   | File Storage   |
| SaaS + Chat   |   | crawler/search |   | uploads        |
+---------------+   +----------------+   +----------------+
\end{diagramblock}

\subsection{اجزای اصلی}
\begin{longtable}{>{\raggedleft\arraybackslash}p{4cm} p{10.5cm}}
\toprule
\textbf{Component} & \textbf{مسئولیت} \\
\midrule
\code{frontend} & رابط مدیریتی برای نقش‌های super admin، customer admin و agent. ارتباط با API از طریق \code{frontend/lib/api.ts}. \\
\code{widget} & ویجت قابل نصب روی سایت مشتری؛ مدیریت visitor، conversation، ارسال پیام، پیوست، نمایش پیام‌ها و درخواست پاسخ AI. \\
\code{backend/api/auth} & ورود، دریافت کاربر جاری، تغییر رمز و خروج از همه دستگاه‌ها. \\
\code{backend/api/widget} & API عمومی با کنترل \code{site\_key}، origin، tenant/site status و rate limit. \\
\code{backend/api/agent} & API عملیات پشتیبانی، مکالمات، پیام‌ها، presence، typing و پیشنهاد AI. \\
\code{backend/api/customer} & مدیریت سایت‌ها، تنظیمات AI، دانش، crawl، team، reports و plan usage. \\
\code{backend/api/super-admin} & مدیریت کل SaaS شامل مشتریان، پلن‌ها، سایت‌ها، کاربران و اعلان‌ها. \\
\code{backend/includes} & هسته مشترک شامل auth، JWT، response، CORS، site access، plan limits، upload، rate limit و AI. \\
\code{MySQL} & نگهداری داده‌های tenant، site، user، visitor، conversation، message، plan، knowledge و AI. \\
\bottomrule
\end{longtable}

\newpage
\section{نمودار استقرار \en{(Deployment Diagram)}}
\textbf{نمودار ۴: Deployment Diagram منطقی}

\begin{diagramblock}
Browser / Visitor
   | loads widget script
   v
Static Widget Host or CDN
   |
   | HTTPS JSON API
   v
PHP Backend Host  <-------------------  Browser / Admin
Apache/Nginx + PHP                      |
backend/api/*                           | HTTPS
   |                                    v
   | PDO                         Next.js Frontend Host
   v                              frontend/app/*
MySQL Server
ai_chat_saas database
   ^
   |
PHP Backend Host ---- file IO ---- Upload Storage
backend/storage/uploads
\end{diagramblock}

\subsection{واحدهای قابل استقرار}
\begin{itemize}
  \item \textbf{Frontend:} پروژه Next.js که می‌تواند روی Node.js server یا سرویس hosting مناسب Next.js اجرا شود.
  \item \textbf{Backend:} فایل‌های PHP زیر وب‌سرور Apache/Nginx با PHP و دسترسی به MySQL.
  \item \textbf{Widget:} فایل جاوااسکریپت static که بهتر است از CDN یا مسیر public پایدار ارائه شود.
  \item \textbf{Database:} MySQL با charset مناسب \code{utf8mb4} برای داده‌های فارسی.
  \item \textbf{Storage:} مسیر ذخیره فایل‌های پیوست و تصاویر اعلان‌ها.
\end{itemize}

\newpage
\section{نمودارهای توالی \en{(Sequence Diagrams)}}
\subsection{جریان ایجاد مشتری توسط سوپرادمین}
\textbf{نمودار ۵: Sequence ایجاد مشتری، سایت و ادمین مشتری}

\begin{diagramblock}
SuperAdmin -> Panel: submit tenant, site, admin, plan
Panel -> customer-create.php: POST payload + Bearer JWT
customer-create.php -> Auth: require_auth + require_role(super_admin)
customer-create.php -> MySQL: validate plan, email, domain
customer-create.php -> MySQL: BEGIN TRANSACTION
customer-create.php -> MySQL: INSERT tenants
customer-create.php -> MySQL: INSERT sites with generated site_key
customer-create.php -> MySQL: INSERT customer_admin user
customer-create.php -> MySQL: INSERT agent_site_access
customer-create.php -> MySQL: COMMIT
customer-create.php -> Panel: tenant_id, site_id, site_key, install_code
\end{diagramblock}

\subsection{جریان مکالمه بازدیدکننده و پاسخ خودکار AI}
\textbf{نمودار ۶: Sequence مکالمه و AI Auto Reply}

\begin{diagramblock}
Visitor -> Widget: open chat and submit first message
Widget -> widget/config.php: GET config by site_key
Widget -> widget/visitor-start.php: POST browser_id, name, phone, email
Widget -> widget/conversation-start.php: POST site_key, visitor_id, source page
conversation-start.php -> Core: validate site_key, origin, active site/tenant
conversation-start.php -> PlanLimits: ensure monthly conversation limit
conversation-start.php -> MySQL: create or reuse active conversation
Widget -> widget/message-send.php: POST visitor message
message-send.php -> MySQL: insert message(sender_type=visitor)
Widget -> widget/ai-reply.php: POST message_id
ai-reply.php -> Core: check ai_mode, settings, duplicate, support online
ai-reply.php -> AI Engine: ai_find_best_answer(site, question)
AI Engine -> MySQL: search knowledge_sources, ai_content_chunks, generated_questions
ai-reply.php -> MySQL: insert AI/fallback message, log, unanswered if needed
Widget -> widget/messages-list.php: poll latest messages
\end{diagramblock}

\subsection{جریان پیشنهاد پاسخ برای Agent}
\textbf{نمودار ۷: Sequence تولید AI Suggestion برای اپراتور}

\begin{diagramblock}
Agent -> Panel: click generate AI suggestion
Panel -> ai-suggestion-generate.php: POST conversation_id + Bearer JWT
ai-suggestion-generate.php -> Auth: require_auth, role customer_admin/agent
ai-suggestion-generate.php -> PlanLimits: require ai_suggestions_enabled
ai-suggestion-generate.php -> SiteAccess: require_site_access
ai-suggestion-generate.php -> MySQL: load conversation and latest visitor message
ai-suggestion-generate.php -> AI Engine: ai_find_best_answer(site, question)
AI Engine -> MySQL: search manual knowledge, chunks, generated questions
ai-suggestion-generate.php -> MySQL: upsert ai_suggestions
ai-suggestion-generate.php -> MySQL: insert ai_answer_logs
ai-suggestion-generate.php -> MySQL: insert unanswered if score is low
ai-suggestion-generate.php -> Panel: suggested_reply, confidence, sources
\end{diagramblock}

\newpage
\section{نمودار دامنه \en{(Class / Domain Diagram)}}
در PHP فعلی، مدل‌ها به‌صورت کلاس‌های domain entity پیاده‌سازی نشده‌اند؛ اما domain model سیستم از جداول و endpointها قابل استخراج است.

\textbf{نمودار ۸: Domain/Class Diagram مفهومی}
\begin{diagramblock}
Plan 1 --- * Tenant
Tenant 1 --- * Site
Tenant 1 --- * User
User * --- * Site       through agent_site_access
Site 1 --- * Visitor
Site 1 --- * Conversation
Visitor 1 --- * Conversation
Conversation 1 --- * Message
Message 1 --- * MessageAttachment
Site 1 --- * KnowledgeSource
Site 1 --- 1 AiSiteSettings
Site 1 --- * AiPage
AiPage 1 --- * AiContentChunk
AiContentChunk 1 --- * AiTerm
AiContentChunk 1 --- * AiGeneratedQuestion
Conversation 1 --- * AiAnswerLog
Conversation 1 --- * AiUnansweredQuestion
\end{diagramblock}

\newpage
\section{نقشه سرویس‌های منطقی \en{(Microservice Diagram)}}
کد فعلی به‌صورت microservice فیزیکی جداگانه پیاده‌سازی نشده است؛ همه APIها در یک بک‌اند PHP قرار دارند. با این حال، معماری داخلی به چند سرویس منطقی تقسیم می‌شود که در آینده می‌توانند به سرویس مستقل تبدیل شوند.

\textbf{نمودار ۹: سرویس‌های منطقی داخل Monolith فعلی}
\begin{diagramblock}
                  +---------------------+
                  | PHP API Monolith    |
                  +----------+----------+
                             |
    +------------------------+------------------------+
    |             |            |          |            |
    v             v            v          v            v
Auth Service  Tenant/Site   Chat/InBox  AI/KM      Reporting
JWT, roles    Plans, users  messages    crawler    dashboards
                           conversations search     usage
    |             |            |          |            |
    +-------------+------------+----------+------------+
                              v
                    Shared MySQL Database
\end{diagramblock}

\newpage
\section{نمودار شبکه \en{(Network Diagram)}}
\textbf{نمودار ۱۰: Network Diagram پیشنهادی برای Production}
\begin{diagramblock}
Internet Browsers
      |
      | HTTPS
      v
Public Domain / TLS Termination
      |
      +----> Next.js Frontend Host
      |
      +----> PHP Backend API Host
      |          |
      |          +---- internal DB connection ----> Private MySQL Server
      |          |
      |          +---- protected file IO --------> Upload Storage
      |
      +----> Static Widget URL / CDN
\end{diagramblock}

\subsection{مرزهای امنیتی شبکه}
\begin{itemize}
  \item دیتابیس نباید مستقیم از اینترنت قابل دسترس باشد و فقط backend مجوز اتصال داشته باشد.
  \item مسیرهای upload باید اجرای اسکریپت را مسدود کنند و فقط فایل‌های مجاز public شوند.
  \item دامنه frontend و backend باید در CORS و origin validation به‌صورت دقیق تنظیم شوند.
  \item ویجت باید از مسیر پایدار HTTPS ارائه شود تا سایت مشتری بتواند script را امن بارگذاری کند.
\end{itemize}

\newpage
\section{نمودار پایگاه داده \en{(Database Diagram)}}
\textbf{نمودار ۱۱: Database Diagram سطح بالا}
\begin{diagramblock}
plans
  |
tenants
  |---------------- users
  |                    |
  |                    +-- agent_site_access -- sites
  |                                             |
  +-------------------- sites -----------------+
                         |
                         +-- visitors
                         |      |
                         |      +-- conversations -- messages -- message_attachments
                         |                 |             |
                         |                 |             +-- ai_answer_logs
                         |                 +-- ai_unanswered_questions
                         |
                         +-- knowledge_sources
                         +-- ai_site_settings
                         +-- ai_crawl_sources -- ai_crawl_runs
                         +-- ai_pages -- ai_content_chunks -- ai_terms
                                      |                  |
                                      |                  +-- ai_generated_questions
                                      +-- ai_answer_logs
\end{diagramblock}

\subsection{توضیح مدل داده}
مدل داده حول محور tenant و site شکل گرفته است. این دو موجودیت مرز اصلی چندمستاجری هستند. کاربران، سایت‌ها، دانش AI، مکالمات، پیام‌ها و گزارش‌ها باید در سطح tenant/site فیلتر شوند. جداول AI در migration موجود با کلیدهای خارجی به \code{tenants} و \code{sites} متصل شده‌اند و برای جست‌وجوی بهتر دارای ایندکس و \en{FULLTEXT KEY} هستند.

\section{معماری جریان داده}
\subsection{جریان داده چت}
\begin{enumerate}
  \item visitor در ویجت، پیام یا فایل ارسال می‌کند.
  \item API عمومی با \code{site\_key} سایت را پیدا می‌کند.
  \item origin و وضعیت tenant/site بررسی می‌شود.
  \item visitor و conversation در دیتابیس ثبت یا بازیابی می‌شوند.
  \item message در جدول \code{messages} ثبت می‌شود.
  \item agent از پنل، پیام‌های conversation را واکشی می‌کند.
  \item پاسخ agent در همان conversation ثبت می‌شود و ویجت در polling بعدی آن را دریافت می‌کند.
\end{enumerate}

\subsection{جریان داده دانش AI}
\begin{enumerate}
  \item customer admin منبع crawl یا دانش دستی تعریف می‌کند.
  \item crawler صفحه را دریافت و محتوای خوانا را استخراج می‌کند.
  \item متن در \code{ai\_pages} ذخیره و به chunk تبدیل می‌شود.
  \item از chunkها terms و generated questions ساخته می‌شود.
  \item هنگام سوال visitor، موتور AI روی knowledge، generated questions و chunks جست‌وجو می‌کند.
  \item پاسخ، confidence و sources در \code{ai\_answer\_logs} ثبت می‌شود.
  \item سوالات کم‌اعتماد وارد \code{ai\_unanswered\_questions} می‌شوند.
\end{enumerate}

\section{تصمیم‌های معماری}
\begin{longtable}{>{\raggedleft\arraybackslash}p{4.2cm} p{6cm} p{4.2cm}}
\toprule
\textbf{تصمیم} & \textbf{دلیل} & \textbf{اثر معماری} \\
\midrule
بک‌اند PHP فایل‌محور & توسعه سریع MVP و endpointهای قابل فهم & سادگی بالا، نیاز به نظم مستندات در رشد آینده \\
چندمستاجری با tenant/site & فروش SaaS به چند مشتری و چند سایت & نیاز دائمی به access control در همه endpointها \\
ویجت با Shadow DOM & جلوگیری از تداخل CSS سایت مشتری & نصب ساده و پایداری UI در سایت‌های مختلف \\
AI استخراجی مبتنی بر دانش & کنترل پاسخ، کاهش hallucination و کاهش وابستگی خارجی & نیاز به crawler، scoring، logs و unanswered loop \\
Plan Limits در backend & مدل درآمدی tiered و کنترل مصرف & enforcement در conversation، agent، site و AI features \\
JWT با token version & احراز هویت stateless همراه با امکان revoke & logout-all و تغییر رمز قابل اعمال می‌شود \\
\bottomrule
\end{longtable}

\section{مقیاس‌پذیری و مسیر تکامل معماری}
معماری فعلی برای MVP و demo فنی مناسب است و مرزبندی ماژول‌ها به‌خوبی قابل مشاهده است. برای رشد تجاری، پیشنهادهای زیر مسیر تکامل طبیعی معماری هستند:
\begin{itemize}
  \item انتقال crawlهای بزرگ از request همزمان به job queue و worker.
  \item افزودن API documentation رسمی برای endpointها.
  \item تفکیک storage عمومی و خصوصی، همراه با signed URL در صورت نیاز.
  \item اضافه کردن cache برای تنظیمات widget و plan usage.
  \item افزودن logging ساختاریافته، monitoring و alerting.
  \item تکمیل migrationهای schema پایه و seedهای demo.
  \item ایجاد تست integration برای جریان‌های tenant/site/access و widget.
\end{itemize}

\section{ریسک‌ها و محدودیت‌های معماری فعلی}
\begin{longtable}{>{\raggedleft\arraybackslash}p{4cm} p{10.5cm}}
\toprule
\textbf{ریسک/محدودیت} & \textbf{تحلیل و اقدام پیشنهادی} \\
\midrule
نبود router مرکزی & endpointهای مستقل ساده‌اند، اما در رشد پروژه باید naming، error handling و includeها استانداردتر شوند. \\
نبود migration کامل schema پایه & برای استقرار و ارزیابی فنی باید schema کامل جداول اصلی کنار migration AI ارائه شود. \\
Crawler همزمان & crawl در request می‌تواند timeout شود؛ worker و queue پیشنهاد می‌شود. \\
SSL verification در crawler & در کد فعلی برای fetch صفحات غیرفعال است؛ در production باید سیاست امن‌تری اعمال شود. \\
AI غیرمولد در جریان فعال & مزیت کنترل‌پذیری دارد، اما باید در ارائه با عنوان موتور دانش استخراجی معرفی شود نه LLM مستقل کامل. \\
تست خودکار محدود & برای بلوغ فنی، تست‌های auth، access، widget flow، upload و AI flow باید اضافه شوند. \\
\bottomrule
\end{longtable}

\section{جمع‌بندی}
معماری \en{AI Chat SaaS} از سه لایه client، API و داده تشکیل شده و برای یک محصول SaaS چندمستاجری طراحی شده است. مرزهای اصلی معماری شامل tenant، site، role، plan و knowledge هستند. ویجت چت، پنل مدیریتی و بک‌اند PHP به‌همراه MySQL و موتور AI استخراجی، یک چرخه کامل پشتیبانی انسانی و هوشمند ایجاد می‌کنند.

نکات قابل تاکید در ارزیابی:
\begin{itemize}
  \item معماری چندمستاجری واقعی و قابل فروش به چند مشتری.
  \item جداسازی نقش‌ها و کنترل دسترسی در سطح سایت.
  \item ترکیب چت انسانی، پاسخ خودکار، پیشنهاد AI و حلقه بهبود دانش.
  \item وجود ساختار plan و limit enforcement برای مدل SaaS.
  \item طراحی ویجت قابل نصب با Shadow DOM و API عمومی کنترل‌شده.
  \item مسیر تکامل روشن برای queue، monitoring، تست، billing و مستندسازی API.
\end{itemize}

\end{document}
